\scanpage{123}
4) \inserted{\handwritten{\uline{\unsure{Trivialite|Trivialité} cohomologique d'}}}\uline{\unsure{Un|un} morphisme projectif et lisse} $f:X\handwritten{\longrightarrow}Y$ \uline{de schémas} (correspondant intuitivement à une fibration sur $Y$ à fibres de variétés projectives non singulières). \deleted{est toujours cohomologiquement trivial (i.e.} \inserted{\handwritten{En d'autres termes,}} \deleted{l}\handwritten{L}a suite spectrale de Leray pour la cohomologie $\handwritten{\unsure{\ell|l}}$-adique - ou à coéffi\deleted{\unsure{s}}cients réels dans le cas du corps de base $C$\deleted{)} - dégénère\deleted{)}. \deleted{Un cas très particulier de ce résultat}\inserted{\handwritten{\deleted{\unsure{E}}Ce théorème est prouvé par Deligne modulo le ``Théorème de \unsure{Lefschetz}'' pour la cohomologie des fibres.}} \handwritten{S}a démonstration s'applique également dans le contexte des variétés analytiques, pour une fibration analytique complexe à fibres compactes kählériennes. Un cas très particulier de ce théorème, relatif au cas où le système local de la cohomologie des fibres était trivial, constitu\deleted{e}\inserted{ait} le résultat principal de la thèse de Blanchard; mais l'hypothèse de Blanchard n'est réalisée dans presque aucun cas intéressant en géométrie. La démonstration de Deligne est différente bien entendu de celle de Blanchard, que Deligne ne connaissait pas.

Déjà du point de vue transcendant, l'intérêt du résultat de Deligne \deleted{\unsure{en top}} pour la topologie des variétés analytiques est évident. Dans l'optique du rédacteur cependant, ce sont ses implications en géométrie algébrique abstraite,et à longue échéance, à la théorie des nombres, qui sont les plus intéressantes\handwritten{,} \deleted{Par exemple, le}\deleted{\unsure{et}} \handwritten{C}e sont d'ailleurs des considérations arithmétiques qui avaient amené le \deleted{rédacteur}\inserted{\handwritten{rapporteur}} à communiquer à Deligne \deleted{la conjecture} l'énoncé précédent \unsure{somme|comme} une conjecture. Grâce au théorème de Deligne, on peut voir par exemple que les hypothèses de \unsure{Reimann|Riemann}-Weil en géométrie diophantienne, interprétées en termes des propriétés d'un endomorphisme de frobénius sur la cohomologie $\handwritten{\ell}$-adique, \deleted{restent valables sous une} impliquent des propriétés analogues dans le cas de coéfficients tordus (pouvant avoir des singularités quelconques, de plus). Le théorème de Deligne sera un des \inserted{\handwritten{nombreux}} ingrédients techniques \deleted{nécessaires}\inserted{\handwritten{indispensables}} pour développer la ``théorie des motifs'', qui \handwritten{\unsure{sera}}
% archive sheet number: 122
